\documentclass[11pt,a4paper]{article}

\usepackage{fontspec}
\usepackage{amssymb}
\usepackage[margin=2.5cm]{geometry}
\usepackage[spanish,es-tabla]{babel}
\usepackage{booktabs}
\usepackage{enumitem}
\usepackage{fancyhdr}
\usepackage{hyperref}
\usepackage{listings}
\usepackage{xcolor}

\hypersetup{colorlinks=true,linkcolor=blue!60!black,urlcolor=blue!60!black}
\setlength{\parskip}{0.5em}
\setlength{\parindent}{0pt}
\setlist{itemsep=0.3em}

\definecolor{sevCritical}{HTML}{B00020}
\definecolor{sevHigh}{HTML}{C75100}
\definecolor{sevMedium}{HTML}{9A6B00}
\definecolor{sevLow}{HTML}{2E7D32}
\definecolor{stOk}{HTML}{1B7F3B}
\definecolor{stPending}{HTML}{B26A00}
\definecolor{stInfo}{HTML}{1565C0}
\definecolor{codebg}{HTML}{F5F5F5}
\definecolor{codenum}{HTML}{9E9E9E}

\newcommand{\badge}[2]{\colorbox{#1}{\textcolor{white}{\strut\footnotesize\bfseries #2}}}
\newcommand{\CRIT}{\badge{sevCritical}{CRÍTICA}}
\newcommand{\HIGH}{\badge{sevHigh}{ALTA}}
\newcommand{\MED}{\badge{sevMedium}{MEDIA}}
\newcommand{\LOW}{\badge{sevLow}{BAJA}}
\newcommand{\FIXED}{\badge{stOk}{CORREGIDO}}
\newcommand{\NREPRO}{\badge{stInfo}{NO REPRODUCE EN FUENTE}}
\newcommand{\MITIG}{\badge{stInfo}{MITIGADO POR DISEÑO}}
\newcommand{\PEND}{\badge{stPending}{PENDIENTE INFRAESTRUCTURA}}
\newcommand{\PVER}{\badge{stPending}{PENDIENTE VERIFICACIÓN}}

\lstdefinestyle{ts}{
  backgroundcolor=\color{codebg},
  basicstyle=\footnotesize\ttfamily,
  breaklines=true,
  breakatwhitespace=true,
  columns=fullflexible,
  commentstyle=\color{green!45!black},
  frame=single,
  framerule=0.4pt,
  keepspaces=true,
  keywordstyle=\color{blue!60!black}\bfseries,
  language=Java,
  morekeywords={string,stringify,readonly,override,async,await},
  numbers=left,
  numbersep=8pt,
  numberstyle=\scriptsize\color{codenum},
  showstringspaces=false,
  stringstyle=\color{red!55!black},
  xleftmargin=1.2em,
}
\lstdefinestyle{sh}{
  backgroundcolor=\color{codebg},
  basicstyle=\footnotesize\ttfamily,
  breaklines=true,
  columns=fullflexible,
  frame=single,
  framerule=0.4pt,
  keepspaces=true,
  language=bash,
  showstringspaces=false,
  xleftmargin=1.2em,
}

\pagestyle{fancy}
\fancyhf{}
\fancyhead[L]{\footnotesize Happy CRM --- Respuesta de remediación}
\fancyhead[R]{\footnotesize Confidencial --- uso interno}
\fancyfoot[C]{\footnotesize \thepage}
\renewcommand{\headrulewidth}{0.4pt}

\begin{document}

\begin{titlepage}
\centering
\vspace*{2cm}
{\LARGE\bfseries Informe de Remediación\par}
\vspace{0.4cm}
{\Large Aplicación CRM (Happy CRM / Twenty)\par}
\vspace{0.6cm}
{\large Respuesta al informe de vulnerabilidades del 26 de octubre de 2026\par}
\vspace{1.2cm}
\begin{tabular}{rl}
\toprule
Para: & Equipo de Seguridad de la Información --- Happy BPO \\
De: & Equipo de desarrollo --- Happy CRM \\
Fecha: & 6 de noviembre de 2026 \\
Referencia: & \texttt{Informe\_Vulnerabilidades\_CRM\_Dev\_2026-10-26.pdf} \\
Código evaluado: & \texttt{main @ 546b8dc665} + rama de trabajo \\
Base upstream: & \texttt{twentyhq/twenty @ ba29ee4f5b} (2026-10-26) + 33 commits \\
Clasificación: & Confidencial --- uso interno / remediación \\
\bottomrule
\end{tabular}
\vfill
{\small Criterio de estados: \textbf{Corregido en fuente} significa cambio implementado y
verificado en código con tests, pendiente de despliegue a desarrollo y retest en runtime.
Ningún estado de este documento sustituye el retest del equipo auditor.\par}
\end{titlepage}

\tableofcontents
\newpage

\section{Resumen Ejecutivo}

El 26 de octubre de 2026 recibimos el informe de vulnerabilidades en caja negra sobre el
entorno de desarrollo, con 6 hallazgos de configuración y despliegue más 5 CVEs de
software base pendientes de verificación por versión. Este documento responde punto por
punto, con evidencia a nivel de código para todo lo que está en manos del equipo de
desarrollo.

El triaje separó el informe en dos grupos. El primero es \textbf{código de la
aplicación}, nuestra responsabilidad directa: aquí corregimos 2 hallazgos, verificamos
que un tercero no reproduce en el fuente actual, confirmamos un cuarto como mitigado por
diseño, y sumamos 2 endurecimientos adicionales fuera del alcance original. El segundo
grupo es \textbf{infraestructura de despliegue} (TLS, puertos publicados, panel
Dokploy): no se puede corregir desde el repositorio, así que dejamos la remediación
exacta especificada, con responsable y criterio de cierre, marcada como bloqueante para
producción.

\begin{center}
\begin{tabular}{lllc}
\toprule
ID & Hallazgo & Severidad & Estado \\
\midrule
DEV-01 & CORS con credenciales & \CRIT & \NREPRO \\
DEV-02 & HTTP plano, puerto 3100 & \HIGH & \PEND \\
DEV-03 & Panel Dokploy expuesto & \MED & \PEND \\
DEV-04 & Certificado Traefik por defecto & \MED & \PEND \\
DEV-05 & Cabecera \texttt{X-Powered-By} & \LOW & \FIXED \\
DEV-06 & Health checks públicos & \LOW & \MITIG \\
\midrule
CVE & 5 CVEs Twenty CRM & Varias & 2 remediados en fuente, 3 por verificar \\
CVE & 5 CVEs Dokploy & Varias & \PVER{} (se requiere versión instalada) \\
\bottomrule
\end{tabular}
\end{center}

Aceptamos la fase 2 autenticada propuesta (IDOR, control multi-tenant y lógica de
negocio) una vez redesplegado el entorno con DEV-01 y DEV-02 atendidos.

\section{Alcance y Metodología}

\textbf{Código evaluado.} Rama \texttt{main} en \texttt{546b8dc665} más la rama de
trabajo con las correcciones de este informe (pendientes de commit a la fecha de
emisión). La base upstream es \texttt{twentyhq/twenty @ ba29ee4f5b} del 2026-10-26
---la misma fecha de la auditoría--- más 33 commits propios, por lo que el análisis
parte de un upstream al día.

\textbf{Método.} Por cada hallazgo: (1) localización del código responsable,
(2) lectura del flujo completo hasta el punto de decisión, (3) revisión de historial
para datar el comportamiento, (4) corrección mínima donde correspondía, y
(5) verificación con la suite del repositorio (jest por archivo, \texttt{oxlint},
\texttt{oxfmt}, \texttt{tsgo --noEmit}). No se probaron endpoints del host auditado:
todo lo afirmado sobre el despliegue se marca explícitamente como pendiente de
verificación en runtime.

\textbf{Fuera de alcance de este documento.} Cambios de infraestructura (se especifican,
no se ejecutan), versiones instaladas en el host (Dokploy, imagen desplegada) y pruebas
autenticadas de lógica de negocio (fase 2).

\section{Remediación detallada}

\subsection[DEV-01 --- CORS con credenciales]{DEV-01 --- CORS con credenciales \hfill \CRIT{} \NREPRO}

\textbf{Conclusión.} El código fuente implementa exactamente la remediación recomendada
(allowlist explícita, sin reflejo dinámico) desde el 2026-08-04. El reflejo observado
en el host auditado no puede provenir de este código: la hipótesis principal es un
build desplegado desactualizado. Se requiere verificar el despliegue y repetir la
prueba (ver acciones al final).

\textbf{Análisis a nivel de código.} El middleware
(\path{user-session/utils/apply-credentialed-cors.util.ts}) refleja el origen solo si
pertenece a la allowlist; cualquier otro recibe literalmente \texttt{*}, que con
\texttt{credentials: true} el navegador bloquea por especificación Fetch:

\begin{lstlisting}[style=ts]
app.enableCors({
  origin: (
    origin: string | undefined,
    callback: (error: Error | null, allow?: boolean | string) => void,
  ) => {
    if (
      origin &&
      resolveAllowedCredentialedOrigins(twentyConfigService).has(
        origin.toLowerCase(),
      )
    ) {
      return callback(null, true);
    }

    return callback(null, '*');
  },
  credentials: true,
\end{lstlisting}

La allowlist se construye exclusivamente desde configuración, sin ningún valor
derivado del request (\path{user-session/utils/resolve-allowed-credentialed-origins.util.ts}):

\begin{lstlisting}[style=ts]
const allowedOriginsSettings: AllowedOriginsSettings = {
  serverUrl: twentyConfigService.get('SERVER_URL'),
  frontendUrl: twentyConfigService.get('FRONTEND_URL'),
  authCookieAllowedOrigins: twentyConfigService.get(
    'AUTH_COOKIE_ALLOWED_ORIGINS',
  ),
  nodeEnvironment: twentyConfigService.get('NODE_ENV'),
};
\end{lstlisting}

Cada candidato debe parsear como URL \texttt{http(s)} (un comodín \texttt{*} se
descarta), los esquemas opacos se rechazan, y en producción los orígenes loopback
derivados se excluyen salvo declaración explícita. El historial confirma este diseño
desde el commit \texttt{997b2c38de} (2026-08-04); la versión anterior al refuerzo
\texttt{51abfb3e40} ya tenía la misma forma. No existe ruta de reflejo dinámico en el
fuente ni en su historial reciente.

\textbf{Discrepancia con el PoC y acciones.} Ordenadas por probabilidad: (1) build
desplegado anterior al 2026-08-04; (2) valores de entorno inesperados en Dokploy;
(3) lectura de \texttt{/healthz} ---público por diseño, donde \texttt{*} es
correcto--- interpretada como lectura autenticada. Nota metodológica: un PoC basado en
curl no prueba lectura cross-origin con credenciales, porque curl no aplica la
política del navegador; solo una prueba en navegador con \texttt{credentials: include}
es concluyente. Acciones: confirmar TAG/commit desplegado, auditar las tres variables
de entorno, redesplegar desde la rama actual y repetir la prueba (comando de retest en
el Anexo~A).

\textbf{Endurecimiento aplicado.} Fallo cerrado: se omite
\texttt{Access-Control-Allow-Origin} para orígenes no permitidos (antes
\texttt{*}). En navegador el observable no cambia (ambos bloquean con
credenciales), pero en runtime ya no existe ninguna cabecera que un escáner
pueda interpretar como aprobación.

\subsection[DEV-02 --- HTTP plano en puerto 3100]{DEV-02 --- HTTP plano en puerto 3100 \hfill \HIGH{} \PEND}

El puerto 3100 es la publicación directa del servidor (\texttt{HOST\_PORT=3100} en
\path{packages/twenty-docker/docker-compose.dokploy.yml}). El propio compose
documenta la remediación pendiente:

\begin{lstlisting}[style=sh]
# En prod con dominio: quitar este publish y rutear solo via Traefik/443.
# Un puerto publicado sigue alcanzable por HTTP plano desde Internet.
ports:
  - "${HOST_PORT:-3000}:3000"
\end{lstlisting}

\textbf{Remediación especificada (bloqueante para producción, owner:
infraestructura):} eliminar el bloque \texttt{ports} del servicio \texttt{server},
servir exclusivamente vía Traefik con TLS válido, redirección a HTTPS y HSTS. No se
aplicó en esta pasada por tratarse de un cambio de despliegue, no de código.

\subsection[DEV-03 --- Panel Dokploy expuesto]{DEV-03 --- Panel Dokploy expuesto \hfill \MED{} \PEND}

Fuera del repositorio de la aplicación. Se aceptan las tres recomendaciones del
informe: restringir por VPN o allowlist de IPs (nunca Internet abierto), confirmar
versión y parchear según los CVEs aplicables, y activar MFA en cuentas
administrativas. Constatamos que las APIs probadas devuelven 401 sin autenticación,
lo cual es correcto pero no sustituye la restricción de red.

\subsection[DEV-04 --- Certificado Traefik por defecto]{DEV-04 --- Certificado Traefik por defecto \hfill \MED{} \PEND}

El certificado \texttt{TRAEFIK DEFAULT CERT} proviene de la configuración TLS de
Dokploy/Traefik, no del código. Acción: asociar el dominio real con certificado
válido (Let's Encrypt vía Dokploy) antes de producción. Bloqueante junto con DEV-02.

\subsection[DEV-05 --- Cabecera \texttt{X-Powered-By}]{DEV-05 --- Cabecera \texttt{X-Powered-By} \hfill \LOW{} \FIXED}

\textbf{Solución aplicada} en \path{packages/twenty-server/src/main.ts} durante el
arranque, antes de registrar el resto de middlewares:

\begin{lstlisting}[style=ts]
// packages/twenty-server/src/engine/utils/disable-powered-by-header.util.ts,
// wired in main.ts and the integration-test harness (create-app.ts)
export const disablePoweredByHeader = (app: INestApplication): void => {
  app.getHttpAdapter().getInstance().disable('x-powered-by');
};
\end{lstlisting}

\textbf{Verificación.} Cambio verificado en fuente; pendiente de commit, despliegue a
desarrollo y retest en runtime (comando en el Anexo~A: la respuesta no debe incluir
la cabecera).

\subsection[DEV-06 --- Health checks públicos]{DEV-06 --- Health checks públicos \hfill \LOW{} \MITIG}

\textbf{Análisis.} El endpoint (\texttt{/healthz}, \texttt{ApiPath.Health}) ejecuta
el chequeo con cero indicadores: devuelve estado \texttt{ok} sin detalle estructurado
de dependencias, versiones o topología:

\begin{lstlisting}[style=ts]
@Get()
@UseGuards(PublicEndpointGuard, NoPermissionGuard)
@HealthCheck()
check() {
  return this.health.check([]);
}
\end{lstlisting}

\textbf{Justificación.} Debe permanecer sin autenticación porque lo consumen el
\texttt{healthcheck} de Docker y las probes del balanceador:

\begin{lstlisting}[style=sh]
healthcheck:
  test: curl --fail http://localhost:3000/healthz
\end{lstlisting}

Quedó registrado en el inventario obligatorio de endpoints públicos con su control
documentado (Sección~5). Su exposición a Internet desaparece automáticamente con la
corrección de DEV-02. Riesgo residual aceptado para desarrollo.

\section{CVEs del software base --- verificación}

\subsection{Determinación de versión}

El fork está al día con \texttt{upstream/main} al 2026-10-26
(\texttt{ba29ee4f5b}, la misma fecha de la auditoría) más 33 commits propios. Los
tags públicos de upstream con esquema \texttt{v}-prefijado llegan hasta
\texttt{v2.9.0} (verificado contra el remoto), por lo que las versiones de
corrección citadas en el informe (\texttt{2.15.0}, \texttt{2.22.0}) no tienen
correspondencia directa visible desde el repositorio. En el despliegue,
\texttt{APP\_VERSION} es el \texttt{TAG} de la imagen Dokploy. Para cerrar el mapeo
se requiere ese TAG más la fuente de versión usada por el informe (Sección~6).

\subsection{Twenty CRM --- evidencia en fuente}

\begin{center}
\begin{tabular}{lp{8.5cm}c}
\toprule
CVE & Evidencia & Estado \\
\midrule
CVE-2026-105763 (9.6) &
La query GraphQL devuelve \texttt{ConnectedAccountPublicDTO}, que omite
\texttt{connectionParameters} crudos y expone solo la forma pública, con
\texttt{password} eliminado por protocolo. &
\FIXED{} \\
CVE-2026-46624 (RCE) &
Sin ningún \texttt{COPY TO} SQL en el código de la aplicación
(búsqueda exhaustiva): la primitiva de RCE descrita no tiene punto de apoyo. &
\FIXED{} \\
CVE-2026-73069 (SQLi) &
Sin verificación puntual en esta pasada. &
\PVER \\
CVE-2026-85055 (permisos de campo) &
Sin verificación puntual; candidato explícito para fase 2 (oráculo de datos). &
\PVER \\
CVE-2026-55583 (IDOR cross-workspace) &
Sin verificación puntual; candidato explícito para fase 2 (multi-tenant). &
\PVER \\
\bottomrule
\end{tabular}
\end{center}

\textbf{Evidencia CVE-2026-105763.} El resolver expone solo el DTO público
(\path{connected-account/resolvers/connected-account.resolver.ts}):

\begin{lstlisting}[style=ts]
@Query(() => [ConnectedAccountPublicDTO])
\end{lstlisting}

cuyo tipo omite los parámetros crudos
(\path{connected-account/dtos/connected-account-public.dto.ts}):

\begin{lstlisting}[style=ts]
export class ConnectedAccountPublicDTO extends OmitType(ConnectedAccountDTO, [
  'connectionParameters',
]) {
\end{lstlisting}

y la construcción de la forma pública elimina secretos
(\path{imap-smtp-caldav-connection/utils/build-public-connection-parameters.util.ts}):

\begin{lstlisting}[style=ts]
const { password: _, ...publicParams } = params;
\end{lstlisting}

\subsection{Dokploy --- pendiente de infraestructura}

Los cinco CVEs (CVE-2026-24841, CVE-2026-72862, CVE-2026-72901, CVE-2026-24840,
CVE-2025-53376/53375) no son verificables desde este repositorio. Acción: obtener la
versión instalada (panel Dokploy o CLI en el host) y actualizar como mínimo a
\texttt{>=0.29.13}. Dada la exposición pública (DEV-03), este punto es prioritario e
independiente del código del CRM.

\section{Endurecimiento adicional}

Fuera del alcance del informe, aprovechamos la pasada para cerrar dos frentes
propios del código, ambos verificados con la suite del repositorio.

\subsection[Oráculo de enumeración de usuarios en login]{Oráculo de enumeración de usuarios en login \hfill \FIXED}

\textbf{Problema.} Email inexistente y password incorrecta devolvían códigos
distintos (\texttt{USER\_NOT\_FOUND} frente a \texttt{FORBIDDEN}), visibles para el
cliente en \texttt{extensions.subCode} de GraphQL: cualquiera podía enumerar cuentas.
\textbf{Solución.} Nuevo punto de paso en
\path{auth/services/auth.service.ts} que mapea el email desconocido al error idéntico
de password incorrecta (mismo código, mensaje y cadena i18n, sin nuevas entradas de
catálogo), usado por las dos rutas de login por password y por la resolución de
workspace en \path{auth/auth.resolver.ts}:

\begin{lstlisting}[style=ts]
// Unknown email maps to the wrong-password error so unauthenticated login
// never reveals whether an email exists. Reuses the exact wrong-password
// strings so no new i18n catalog entry is created.
async findUserByEmailForPasswordLoginOrThrow(email: string) {
  try {
    return await this.findUserByEmailOrThrow(email);
  } catch (error) {
    if (
      error instanceof AuthException &&
      error.code === AuthExceptionCode.USER_NOT_FOUND
    ) {
      throw new AuthException(
        'Wrong password',
        AuthExceptionCode.FORBIDDEN_EXCEPTION,
        {
          userFriendlyMessage: msg`Wrong password.`,
        },
      );
    }

    throw error;
  }
}
\end{lstlisting}

No se tocó la utilidad compartida de validación de password: en el cambio de
password autenticado el mensaje específico sigue siendo el UX correcto. Se verificó
que el front no depende de ninguno de los dos códigos.

\textbf{Verificación.} 2 tests de regresión nuevos; suites
\texttt{auth.service.spec} (7 tests) y \texttt{auth.resolver.spec} (11 tests) en
verde, más \texttt{oxlint}, \texttt{oxfmt} y \texttt{tsgo --noEmit} limpios.

\subsection[Inventario obligatorio de endpoints públicos]{Inventario obligatorio de endpoints públicos \hfill \FIXED}

\textbf{Problema.} \texttt{PublicEndpointGuard} es un marcador que siempre pasa: cada
endpoint público depende de su propio control compensatorio (captcha, firma, token,
throttle), y nada impedía añadir superficie pública sin ninguno.
\textbf{Solución.} Nuevo test
(\path{engine/guards/__tests__/public-endpoint-inventory.spec.ts}) que recorre el
fuente, detecta todos los usos del guard y falla si alguno no está registrado en la
allowlist con su control documentado (o si sobra una entrada):

\begin{lstlisting}[style=ts]
// PublicEndpointGuard always passes, so every file using it must document its
// compensating control here. Add new public surface to the allowlist with the
// control that protects it, or the test fails.
\end{lstlisting}

\begin{lstlisting}[style=ts]
expect(unregistered).toEqual([]);
expect(stale).toEqual([]);
\end{lstlisting}

\textbf{Verificación.} Superficie actual: 32 archivos, 16 con control verificado en
esta pasada (incluidos webhooks con firma, throttle en 2FA y tracking, tokens de
256 bits y archivos públicos con doble validación de ruta); el resto queda marcado
como pendiente de revisión en el propio test, con revisión sugerida de
\texttt{workflow-trigger.controller.ts} y el OAuth de applications.

\subsection{Verificado sin cambios necesarios}

Throttle existente en 2FA y captura de tracking, tokens de verificación de 256 bits
con expiración y hash en reposo, firma de webhooks (Stripe sobre \texttt{rawBody};
Svix/SNS delegadas a drivers), archivos públicos con doble validación de ruta más
registro en base de datos, y \texttt{isAuthRequired} con default público explícito de
producto en rutas HTTP (decisión de UX tipo webhook, documentada).

\section{Trabajo pendiente y compromisos}

Lo que está en manos del equipo de desarrollo quedó corregido y verificado en fuente.
Lo restante es infraestructura y verificación conjunta:

\begin{itemize}
\item[$\square$] Commit + despliegue a desarrollo de DEV-05 y los endurecimientos
de la Sección~5, con retest en runtime (Anexo~A). Owner: desarrollo.
\item[$\boxtimes$] Endurecimiento CORS fail-closed, pendiente de despliegue y
retest. Owner: desarrollo.
\item[$\boxtimes$] Neutralización de branding extendida a
\texttt{getPublicWorkspaceDataById} (misma clase de leak en DNS compartido),
pendiente de despliegue y retest. Owner: desarrollo.
\item[$\square$] Cierre del puerto 3100 público + TLS válido + restricción de
Dokploy \textbf{antes de producción} (bloqueantes). Owner: infraestructura.
\item[$\square$] Fase 2 autenticada con usuario de prueba (cubre CVE-2026-73069,
CVE-2026-85055 y CVE-2026-55583 pendientes). Owner: auditoría + desarrollo.
\end{itemize}

\textbf{Información solicitada para cerrar la verificación:}
\begin{enumerate}
\item TAG (o commit) de la imagen \texttt{happy/crm\_happy} desplegada en el host
auditado el 2026-10-26.
\item Valores de \texttt{SERVER\_URL}, \texttt{FRONTEND\_URL} y
\texttt{AUTH\_COOKIE\_ALLOWED\_ORIGINS} del environment Dokploy.
\item Versión instalada de Dokploy.
\item Confirmación y credenciales de prueba para la fase 2.
\end{enumerate}

\appendix
\section{Comandos de retest}

\begin{lstlisting}[style=sh]
# DEV-05: la respuesta no debe incluir la cabecera
curl -sI http://HOST:3100/healthz | grep -i "x-powered-by"
# Esperado: sin salida

# DEV-01: origen no permitido no debe reflejarse
curl -s -o /dev/null -D - -H "Origin: https://evil.example" \
  http://HOST:3100/graphql | grep -i "access-control-allow-origin"
# Esperado: sin salida (ninguna cabecera ACAO = denegado)
# Y la prueba en navegador con credentials:include, que debe bloquearse

# DEV-06: health sin detalle estructurado
curl -s http://HOST:3100/healthz
# Esperado: solo estado ok, sin indicadores ni versiones
\end{lstlisting}

\section{Endpoints públicos con control verificado}

Superficie con control verificado en esta pasada (15 de 32 archivos; el resto figura
como pendiente de revisión en el propio test de inventario):

\begin{center}
\footnotesize
\begin{tabular}{p{6.2cm}p{8.3cm}}
\toprule
Archivo & Control verificado \\
\midrule
\texttt{messaging-webhooks.controller.ts} & Firmas de provider verificadas en drivers SNS/Svix \\
\texttt{unsubscribe.controller.ts} & Throttle por solicitante y por token \\
\texttt{campaign-tracking.controller.ts} & Throttle en el servicio de captura \\
\texttt{connected-account-sync-webhooks.controller.ts} & Tokens de canal validados en handlers \\
\texttt{route-trigger.controller.ts} & Scope por workspace, auth opcional por ruta, rate limit \\
\texttt{email-verification.resolver.ts} & Tokens de 256 bits, expirables, de un uso \\
\texttt{billing-webhook.controller.ts} & Firma Stripe sobre cuerpo crudo \\
\texttt{two-factor-authentication.resolver.ts} & Login-token + throttle de OTP en servicio \\
\texttt{file.controller.ts} & Assets públicos por diseño; doble validación de ruta \\
\texttt{auth.resolver.ts} & CaptchaGuard en credenciales; resto con token \\
\texttt{google-auth.controller.ts} & Flujo OAuth con guard de provider \\
\texttt{sso-auth.controller.ts} & Guards SAML/OIDC; metadata pública por diseño \\
\texttt{google-apis-auth.controller.ts} & Guards OAuth + tokens transitorios \\
\texttt{server-route-trigger.controller.ts} & Despacho público; rechaza resolvers con auth \\
\texttt{health.controller.ts} & Chequeo sin indicadores; requerido por Docker \\
\bottomrule
\end{tabular}
\end{center}

\end{document}
